% Options for packages loaded elsewhere
\PassOptionsToPackage{unicode}{hyperref}
\PassOptionsToPackage{hyphens}{url}
\documentclass[
  mongolian,
]{article}
\usepackage{xcolor}
\usepackage{amsmath,amssymb}
\setcounter{secnumdepth}{-\maxdimen} % remove section numbering
\usepackage{iftex}
\ifPDFTeX
  \usepackage[T1]{fontenc}
  \usepackage[utf8]{inputenc}
  \usepackage{textcomp} % provide euro and other symbols
\else % if luatex or xetex
  \usepackage{unicode-math} % this also loads fontspec
  \defaultfontfeatures{Scale=MatchLowercase}
  \defaultfontfeatures[\rmfamily]{Ligatures=TeX,Scale=1}
\fi
\usepackage{lmodern}
\ifPDFTeX\else
  % xetex/luatex font selection
\fi
% Use upquote if available, for straight quotes in verbatim environments
\IfFileExists{upquote.sty}{\usepackage{upquote}}{}
\IfFileExists{microtype.sty}{% use microtype if available
  \usepackage[]{microtype}
  \UseMicrotypeSet[protrusion]{basicmath} % disable protrusion for tt fonts
}{}
\makeatletter
\@ifundefined{KOMAClassName}{% if non-KOMA class
  \IfFileExists{parskip.sty}{%
    \usepackage{parskip}
  }{% else
    \setlength{\parindent}{0pt}
    \setlength{\parskip}{6pt plus 2pt minus 1pt}}
}{% if KOMA class
  \KOMAoptions{parskip=half}}
\makeatother
\usepackage{longtable,booktabs,array}
\usepackage{calc} % for calculating minipage widths
% Correct order of tables after \paragraph or \subparagraph
\usepackage{etoolbox}
\makeatletter
\patchcmd\longtable{\par}{\if@noskipsec\mbox{}\fi\par}{}{}
\makeatother
% Allow footnotes in longtable head/foot
\IfFileExists{footnotehyper.sty}{\usepackage{footnotehyper}}{\usepackage{footnote}}
\makesavenoteenv{longtable}
\usepackage{graphicx}
\makeatletter
\newsavebox\pandoc@box
\newcommand*\pandocbounded[1]{% scales image to fit in text height/width
  \sbox\pandoc@box{#1}%
  \Gscale@div\@tempa{\textheight}{\dimexpr\ht\pandoc@box+\dp\pandoc@box\relax}%
  \Gscale@div\@tempb{\linewidth}{\wd\pandoc@box}%
  \ifdim\@tempb\p@<\@tempa\p@\let\@tempa\@tempb\fi% select the smaller of both
  \ifdim\@tempa\p@<\p@\scalebox{\@tempa}{\usebox\pandoc@box}%
  \else\usebox{\pandoc@box}%
  \fi%
}
% Set default figure placement to htbp
\def\fps@figure{htbp}
\makeatother
\ifLuaTeX
\usepackage[bidi=basic,shorthands=off,,provide=*]{babel}
\else
\usepackage[bidi=default,shorthands=off,,provide=*]{babel}
\fi
\ifLuaTeX
  \usepackage{selnolig} % disable illegal ligatures
\fi
\setlength{\emergencystretch}{3em} % prevent overfull lines
\providecommand{\tightlist}{%
  \setlength{\itemsep}{0pt}\setlength{\parskip}{0pt}}
\usepackage{fontspec}
\setmainfont{Liberation Serif}
\setsansfont{Liberation Sans}
\setmonofont{Liberation Mono}[Scale=0.86]
\usepackage[a4paper,left=3cm,right=2cm,top=2cm,bottom=2cm,includefoot,headheight=15pt]{geometry}
\usepackage{ragged2e,indentfirst,xurl}
\usepackage{longtable,booktabs,array,calc,graphicx}
\usepackage{xcolor,fancyhdr,titlesec,enumitem,float}
\definecolor{ink}{HTML}{19354A}
\definecolor{accent}{HTML}{167D8D}
\titleformat{\section}{\Large\sffamily\bfseries\color{ink}}{}{0pt}{}
\titleformat{\subsection}{\large\sffamily\bfseries\color{ink}}{}{0pt}{}
\titleformat{\subsubsection}{\fontsize{11.5}{13.5}\selectfont\sffamily\bfseries\color{ink}}{}{0pt}{}
\titlespacing*{\section}{1.25cm}{0pt}{8pt}
\titlespacing*{\subsection}{1.25cm}{8pt}{4pt}
\titlespacing*{\subsubsection}{1.25cm}{6pt}{3pt}
\setlength{\parindent}{1.25cm}
\setlength{\parskip}{3pt}
\setlength{\emergencystretch}{3em}
\setlength{\tabcolsep}{4pt}
\renewcommand{\arraystretch}{1.08}
\setlist{nosep,leftmargin=1.5em,topsep=3pt}
\pagestyle{fancy}\fancyhf{}
\fancyhead[L]{\small\sffamily ICSI438 / SEMINAR 04}
\fancyhead[R]{\small\sffamily Nogoolin / M4}
\fancyfoot[R]{\small\thepage}
\renewcommand{\headrulewidth}{0.3pt}
\AtBeginDocument{\fontsize{10.5}{12.5}\selectfont\justifying\setlength{\parindent}{1.25cm}}


\AtBeginDocument{\renewcommand{\figurename}{Зураг}}
\usepackage{bookmark}
\IfFileExists{xurl.sty}{\usepackage{xurl}}{} % add URL line breaks if available
\urlstyle{same}
\hypersetup{
  pdflang={mn},
  hidelinks,
  pdfcreator={LaTeX via pandoc}}

\author{}
\date{}

\begin{document}

\hypersetup{
  pdftitle={Програм хангамжийн төслийн баримт бичиг — Seminar 4 Architecture},
  pdfauthor={Ц. Тэнгис}
}
\pagenumbering{gobble}
\thispagestyle{empty}
\noindent
\begin{minipage}[t][\dimexpr\textheight-\topskip\relax][t]{\linewidth}
  \begin{center}
    \textbf{МОНГОЛ УЛСЫН ИХ СУРГУУЛЬ}\\
    \textbf{МЭДЭЭЛЛИЙН ТЕХНОЛОГИ, ЭЛЕКТРОНИКИЙН СУРГУУЛЬ}\\
    \textbf{МЭДЭЭЛЭЛ, КОМПЬЮТЕРЫН УХААНЫ ТЭНХИМ}\\[1.7cm]
    \includegraphics[width=6cm]{../assets/num-logo.jpg}\\[1.2cm]
    {\fontsize{16}{20}\selectfont\bfseries
      ПРОГРАМ ХАНГАМЖИЙН ТӨСЛИЙН БАРИМТ БИЧИГ\par}
    \vspace{0.45cm}
    {\large\bfseries SEMINAR 4 / M4 — SOFTWARE ARCHITECTURE DOCUMENTATION}
  \end{center}

  \vspace{1.2cm}
  \begin{tabular}{>{\raggedright\arraybackslash}m{5cm}@{\hspace{1cm}}>{\raggedright\arraybackslash}m{\dimexpr\linewidth-6cm-2\tabcolsep\relax}}
    Хичээл: & ICSI438 — Software Project Documentation \\[0.55cm]
    Үндсэн төсөл: & Nogoolin \\[0.55cm]
    Хичээлийн багш: & Ph.D. Batnyam Battulga \\[0.55cm]
    Гүйцэтгэсэн: & Програм хангамжийн IV ангийн оюутан\\
                  & Ц. Тэнгис (22B1NUM6249) \\[0.55cm]
    Баримтын хувилбар: & v0.1 draft
  \end{tabular}

  \vfill
  \begin{center}
    Улаанбаатар\\[0.25cm]
    2026 он
  \end{center}
\end{minipage}
\clearpage
\pagenumbering{arabic}
\setcounter{page}{1}



\section{Nogoolin архитектурын баримт
бичиг}\label{nogoolin-ux430ux440ux445ux438ux442ux435ux43aux442ux443ux440ux44bux43d-ux431ux430ux440ux438ux43cux442-ux431ux438ux447ux438ux433}

\textbf{arc42 v9 lean profile · ICSI438 Seminar 4 / M4 · Тэнгис · v0.1
draft · 2026-09-28}

Энэ баримт нь Nogoolin-ийн одоогийн код, Seminar 3-ын SRS болон
шийдвэрийн бүртгэлийг нэг архитектурын зураглалд холбов. Arc42-ийн 12
хэсгийг хадгалж, §3, §5, §6-д C4 загварыг хэрэглэв. Week 04 Lab
Assignment-ийн §1, §3, §5, гурван MADR, persona холбоо, ≥80\% FR
coverage болон UE-4 шаардлагыг энэ хувилбарт тусгав. Lecture 04 хараахан
байхгүй.

\begin{longtable}[]{@{}
  >{\raggedright\arraybackslash}p{(\linewidth - 2\tabcolsep) * \real{0.5000}}
  >{\raggedright\arraybackslash}p{(\linewidth - 2\tabcolsep) * \real{0.5000}}@{}}
\toprule\noalign{}
\begin{minipage}[b]{\linewidth}\raggedright
Талбар
\end{minipage} & \begin{minipage}[b]{\linewidth}\raggedright
Утга
\end{minipage} \\
\midrule\noalign{}
\endhead
\bottomrule\noalign{}
\endlastfoot
Төсөл & Nogoolin \\
Зохиогч & Тэнгис \\
Хувилбар & v0.1 draft \\
Requirements baseline &
\href{https://github.com/Tengis01/Nogoolin/blob/main/docs/requirements/srs-v1.0.md}{SRS
v1.0 draft},
\href{https://github.com/Tengis01/Nogoolin/blob/main/docs/requirements/traceability-matrix-v1.0.md}{20
мөртэй matrix} \\
Public repository & \url{https://github.com/Tengis01/Nogoolin} \\
Аргачлал & \href{https://arc42.org/download/}{arc42 v9},
\href{https://c4model.com/}{C4 model} ба ADR \\
\end{longtable}

\section{1. Оршил ба
зорилго}\label{ux43eux440ux448ux438ux43b-ux431ux430-ux437ux43eux440ux438ux43bux433ux43e}

\subsection{1.1 Системийн
зорилго}\label{ux441ux438ux441ux442ux435ux43cux438ux439ux43d-ux437ux43eux440ux438ux43bux433ux43e}

Nogoolin нь Монголын шашны бүтээгдэхүүний каталог, хэрэглэгчийн inquiry,
хадгалсан бүтээгдэхүүний жагсаалт болон админ удирдлагыг web, mobile,
REST API-аар нийлүүлэх систем. Одоогийн M4 хүрээ нь SRS-д тогтоосон
inquiry, wishlist, admin inquiry болон нүүрний fallback урсгалд
төвлөрнө.

\subsection{1.2 Оролцогч
талууд}\label{ux43eux440ux43eux43bux446ux43eux433ux447-ux442ux430ux43bux443ux443ux434}

\begin{longtable}[]{@{}
  >{\raggedright\arraybackslash}p{(\linewidth - 2\tabcolsep) * \real{0.5000}}
  >{\raggedright\arraybackslash}p{(\linewidth - 2\tabcolsep) * \real{0.5000}}@{}}
\toprule\noalign{}
\begin{minipage}[b]{\linewidth}\raggedright
Оролцогч
\end{minipage} & \begin{minipage}[b]{\linewidth}\raggedright
Архитектураас хүлээх зүйл
\end{minipage} \\
\midrule\noalign{}
\endhead
\bottomrule\noalign{}
\endlastfoot
Зочин ба хэрэглэгч & Каталог үзэх, inquiry илгээх, өөрийн өгөгдөлд л
хандах \\
Админ & Inquiry-г шүүх, төлөв солих, өөрчлөлтийг audit-д үлдээх \\
Тэнгис, хөгжүүлэгч & Нэг repository-д ойлгомжтой бүтэц, локал шалгалт,
бага operational overhead \\
ICSI438 багш & SRS → architecture → decision холбоог шалгах боломж \\
\end{longtable}

\subsection{1.3 Архитектурын гол
зорилгууд}\label{ux430ux440ux445ux438ux442ux435ux43aux442ux443ux440ux44bux43d-ux433ux43eux43b-ux437ux43eux440ux438ux43bux433ux443ux443ux434}

\begin{enumerate}
\def\labelenumi{\arabic{enumi}.}
\tightlist
\item
  \textbf{Өгөгдөл тусгаарлалт:} JWT, ownership filter, RBAC, RLS-ийг
  давхар хэрэглэж NFR-01-ийг хангах.
\item
  \textbf{Засварлах боломж:} Controller → Service → Repository хилийг
  хадгалж, UI/API/schema drift-ийг бууруулах.
\item
  \textbf{Хүртээмж ба fallback:} хөдөлгөөн багасгах тохиргоо, static
  hero-г хадгалж NFR-02-ыг хангах.
\item
  \textbf{Шалгаж болох build:} integration test үнэхээр ажилласан
  эсэхийг CI ялгаж NFR-03-ын false-green эрсдэлийг арилгах.
\item
  \textbf{Хариу өгөх хугацаа:} каталогийн замыг NFR-04-ийн 2.0 секундийн
  босготой уялдуулах.
\end{enumerate}

\section{2.
Хязгаарлалт}\label{ux445ux44fux437ux433ux430ux430ux440ux43bux430ux43bux442}

\begin{longtable}[]{@{}
  >{\raggedright\arraybackslash}p{(\linewidth - 2\tabcolsep) * \real{0.5000}}
  >{\raggedright\arraybackslash}p{(\linewidth - 2\tabcolsep) * \real{0.5000}}@{}}
\toprule\noalign{}
\begin{minipage}[b]{\linewidth}\raggedright
Төрөл
\end{minipage} & \begin{minipage}[b]{\linewidth}\raggedright
Хязгаарлалт
\end{minipage} \\
\midrule\noalign{}
\endhead
\bottomrule\noalign{}
\endlastfoot
Баг & Нэг хөгжүүлэгч; microservice-ийн нэмэлт ажиллагааг даахгүй \\
Технологи & TypeScript monorepo; Next.js, Expo, Fastify, Supabase,
Zod \\
Архитектур & Нэг Fastify backend; controller/service/repository
давхарга \\
Өгөгдөл & Local-first Supabase CLI; cloud project Phase 6 хүртэл
үүсээгүй \\
Аюулгүй байдал & Secret зөвхөн environment variable-д; RLS бүх
хүснэгтэд \\
Хүргэлт & \texttt{delivery\_enabled=false}; order урсгал M4-ийн идэвхтэй
scope биш \\
Нотолгоо & Planned test-ийг executed/verified гэж тэмдэглэхгүй \\
\end{longtable}

\section{3. Контекст ба хамрах
хүрээ}\label{ux43aux43eux43dux442ux435ux43aux441ux442-ux431ux430-ux445ux430ux43cux440ux430ux445-ux445ux4afux440ux44dux44d}

Arc42 §3 нь системийн хил, хэрэглэгч болон хөрш системийг харуулна. C4
System Context нь Nogoolin-ийг нэг black box болгон үзүүлж, дотоод
технологийн дэлгэрэнгүйг §5 руу үлдээв.

\begin{figure}[H]
\centering
\includegraphics[width=1\linewidth,height=0.72\textheight,keepaspectratio,alt={C4-01 --- Nogoolin System Context}]{../tmp/diagrams/c4-01-system-context.pdf}
\caption{C4-01 --- Nogoolin System Context}
\end{figure}

\emph{Mermaid эх:
\href{https://github.com/Tengis01/Nogoolin/blob/main/docs/ICSI438/sem4/diagrams/c4-01-system-context.mmd}{c4-01-system-context.mmd}.}

\textbf{Зургийн тэмдэглэгээ:} Person нь хүний role, Software System нь
нэг систем, Container нь ажиллах application, Component нь API доторх
нэгж. Цэнхэр хүрээ Nogoolin-ийн эзэмшил, саарал хүрээ external system.
Сум бүр data flow label-тай. Mermaid flowchart дээр C4 type, name,
technology, responsibility-г ил тод бичсэн; runtime нь C4-ийн зөвшөөрсөн
sequence style.

\subsection{3.1 Бизнес
контекст}\label{ux431ux438ux437ux43dux435ux441-ux43aux43eux43dux442ux435ux43aux441ux442}

\begin{longtable}[]{@{}
  >{\raggedright\arraybackslash}p{(\linewidth - 4\tabcolsep) * \real{0.3333}}
  >{\raggedright\arraybackslash}p{(\linewidth - 4\tabcolsep) * \real{0.3333}}
  >{\raggedright\arraybackslash}p{(\linewidth - 4\tabcolsep) * \real{0.3333}}@{}}
\toprule\noalign{}
\begin{minipage}[b]{\linewidth}\raggedright
Харилцагч
\end{minipage} & \begin{minipage}[b]{\linewidth}\raggedright
Nogoolin руу орох мэдээлэл
\end{minipage} & \begin{minipage}[b]{\linewidth}\raggedright
Nogoolin-оос гарах мэдээлэл
\end{minipage} \\
\midrule\noalign{}
\endhead
\bottomrule\noalign{}
\endlastfoot
Зочин & Каталогийн хайлт, inquiry form & Бүтээгдэхүүн, inquiry
confirmation \\
Хэрэглэгч & Session, inquiry, wishlist үйлдэл & Өөрийн inquiry,
хадгалсан бүтээгдэхүүн \\
Админ & Inquiry filter, status command & Inbox, шинэ төлөв, audit үр
дүн \\
OAuth provider & Authorization response & OAuth authorization request \\
Supabase platform & Auth/storage/database response & Auth, query,
storage request \\
\end{longtable}

\subsection{3.2 Техникийн
контекст}\label{ux442ux435ux445ux43dux438ux43aux438ux439ux43d-ux43aux43eux43dux442ux435ux43aux441ux442}

\begin{itemize}
\tightlist
\item
  Browser ба mobile client нь HTTPS JSON API-аар Fastify backend-тай
  харилцана.
\item
  Web/mobile нь Supabase Auth-тай session үүсгэнэ; API JWT-г шалгана.
\item
  Repository давхарга PostgreSQL болон Storage-тай Supabase client-аар
  харилцана.
\item
  Production-ийн Cloudflare, Vercel, Railway нь төлөвлөсөн target; local
  development-д localhost ба Supabase CLI stack хэрэглэгдэнэ.
\end{itemize}

\subsection{3.3 W1 persona-тай
холбоо}\label{w1-persona-ux442ux430ux439-ux445ux43eux43bux431ux43eux43e}

W1-д нэг хөгжүүлэгч persona, Тэнгис, болон P1/P2/P3 pain point байсан.
Тэдгээрийг шинэ хэрэглэгчтэй хийсэн ярилцлага гэж өөрчлөхгүй. Доорх
гурван external role нь Nogoolin-ийн бизнес actor; W1 persona-ийн
оношлох хэрэгцээг role бүрийн architecture boundary-д шилжүүлэн холбов.

\begin{longtable}[]{@{}
  >{\raggedright\arraybackslash}p{(\linewidth - 4\tabcolsep) * \real{0.3333}}
  >{\raggedright\arraybackslash}p{(\linewidth - 4\tabcolsep) * \real{0.3333}}
  >{\raggedright\arraybackslash}p{(\linewidth - 4\tabcolsep) * \real{0.3333}}@{}}
\toprule\noalign{}
\begin{minipage}[b]{\linewidth}\raggedright
External actor
\end{minipage} & \begin{minipage}[b]{\linewidth}\raggedright
W1 холбоо
\end{minipage} & \begin{minipage}[b]{\linewidth}\raggedright
Nogoolin-д тайлбарласан холбоо
\end{minipage} \\
\midrule\noalign{}
\endhead
\bottomrule\noalign{}
\endlastfoot
Зочин &
\href{https://github.com/Tengis01/Nogoolin/blob/main/docs/ICSI438/sem1/wiki/02_audience_persona.md}{P3:
local/pipeline ялгаа} & Inquiry response ба алдаа аль API/DB хэсэгт
гарсныг ялгах \\
Customer &
\href{https://github.com/Tengis01/Nogoolin/blob/main/docs/ICSI438/sem1/wiki/02_audience_persona.md}{P1:
session алдагдах} & Session, own inquiry, wishlist boundary; P-NG-01 \\
Админ &
\href{https://github.com/Tengis01/Nogoolin/blob/main/docs/ICSI438/sem1/wiki/02_audience_persona.md}{P1,
P3} & JWT/role болон data query/audit-ийг тусад нь шалгах \\
\end{longtable}

Энэ нь гурван role-ийг W1-ийн нэг persona-тай холбосон mapping; W1-д
гурван тусдаа persona байсан гэж зарлаагүй.
\href{https://github.com/Tengis01/Nogoolin/blob/main/docs/ICSI438/sem2/wiki/persona-and-pivot.md}{W2
pivot} нь шилжилтийн үндэслэлийг хадгална.

\section{4. Шийдлийн
стратеги}\label{ux448ux438ux439ux434ux43bux438ux439ux43d-ux441ux442ux440ux430ux442ux435ux433ux438}

\begin{longtable}[]{@{}
  >{\raggedright\arraybackslash}p{(\linewidth - 4\tabcolsep) * \real{0.3333}}
  >{\raggedright\arraybackslash}p{(\linewidth - 4\tabcolsep) * \real{0.3333}}
  >{\raggedright\arraybackslash}p{(\linewidth - 4\tabcolsep) * \real{0.3333}}@{}}
\toprule\noalign{}
\begin{minipage}[b]{\linewidth}\raggedright
Зорилго/хязгаарлалт
\end{minipage} & \begin{minipage}[b]{\linewidth}\raggedright
Стратеги
\end{minipage} & \begin{minipage}[b]{\linewidth}\raggedright
Холбогдох шийдвэр
\end{minipage} \\
\midrule\noalign{}
\endhead
\bottomrule\noalign{}
\endlastfoot
Solo developer, бага ажиллагаа & Нэг Fastify backend бүхий layered
monolith &
\href{https://github.com/Tengis01/Nogoolin/blob/main/docs/architecture/adr/adr-007-layered-monolith.md}{ADR-007} \\
TypeScript REST API & Fastify plugin/hook болон schema validation &
\href{https://github.com/Tengis01/Nogoolin/blob/main/docs/architecture/adr/adr-003-fastify.md}{ADR-003} \\
Relational data, auth, storage, RLS & Supabase platform, repository
abstraction &
\href{https://github.com/Tengis01/Nogoolin/blob/main/docs/architecture/adr/adr-005-supabase.md}{ADR-005} \\
SEO ба admin/public web & Next.js App Router; одоогийн ADR-006 & \\
Shared contract & Zod schema-г \texttt{packages/validation-schemas}-д
төвлөрүүлэх & \\
Local-first security test & Supabase CLI migrations; одоогийн ADR-011
& \\
\end{longtable}

Энэ хэсэг сонголтыг товч харуулна. Static бүтэц §5, runtime харилцан
үйлчлэл §6, бүрэн decision reasoning §9-ийн ADR холбоосуудад байна.

\section{5. Бүрэлдэхүүний
зураглал}\label{ux431ux4afux440ux44dux43bux434ux44dux445ux4afux4afux43dux438ux439-ux437ux443ux440ux430ux433ux43bux430ux43b}

\subsection{5.1 Level 1 --- Container
view}\label{level-1-container-view}

\begin{figure}[H]
\centering
\includegraphics[width=1\linewidth,height=0.72\textheight,keepaspectratio,alt={C4-02 --- Nogoolin Container View}]{../tmp/diagrams/c4-02-container.pdf}
\caption{C4-02 --- Nogoolin Container View}
\end{figure}

\emph{Mermaid эх:
\href{https://github.com/Tengis01/Nogoolin/blob/main/docs/ICSI438/sem4/diagrams/c4-02-container.mmd}{c4-02-container.mmd}.}

\begin{longtable}[]{@{}
  >{\raggedright\arraybackslash}p{(\linewidth - 4\tabcolsep) * \real{0.3333}}
  >{\raggedright\arraybackslash}p{(\linewidth - 4\tabcolsep) * \real{0.3333}}
  >{\raggedright\arraybackslash}p{(\linewidth - 4\tabcolsep) * \real{0.3333}}@{}}
\toprule\noalign{}
\begin{minipage}[b]{\linewidth}\raggedright
Container
\end{minipage} & \begin{minipage}[b]{\linewidth}\raggedright
Үүрэг
\end{minipage} & \begin{minipage}[b]{\linewidth}\raggedright
Технологи ба code location
\end{minipage} \\
\midrule\noalign{}
\endhead
\bottomrule\noalign{}
\endlastfoot
Web + Admin & Public catalog, inquiry, wishlist, profile, admin inbox &
Next.js; \texttt{apps/web} \\
Mobile & Catalog, auth; inquiry/wishlist mobile хэсэг дутуу & Expo/React
Native; \texttt{apps/mobile} \\
REST API & Auth boundary, business rule, persistence orchestration &
Fastify/TypeScript; \texttt{backend/api} \\
Shared validation library & Web/API contract; deployable container биш &
Zod; \texttt{packages/validation-schemas} \\
Supabase Auth & Session, JWT, OAuth bridge & Local stack; cloud
deferred \\
PostgreSQL & Product, inquiry, cart, audit ба RLS &
\texttt{supabase/migrations} \\
Storage & Product media & Supabase Storage \\
\end{longtable}

\subsection{5.2 Level 2 --- Fastify API component
view}\label{level-2-fastify-api-component-view}

\begin{figure}[H]
\centering
\includegraphics[width=1\linewidth,height=0.72\textheight,keepaspectratio,alt={C4-03 --- Fastify API Component View}]{../tmp/diagrams/c4-03-api-component.pdf}
\caption{C4-03 --- Fastify API Component View}
\end{figure}

\emph{Mermaid эх:
\href{https://github.com/Tengis01/Nogoolin/blob/main/docs/ICSI438/sem4/diagrams/c4-03-api-component.mmd}{c4-03-api-component.mmd}.}

\begin{longtable}[]{@{}
  >{\raggedright\arraybackslash}p{(\linewidth - 4\tabcolsep) * \real{0.3333}}
  >{\raggedright\arraybackslash}p{(\linewidth - 4\tabcolsep) * \real{0.3333}}
  >{\raggedright\arraybackslash}p{(\linewidth - 4\tabcolsep) * \real{0.3333}}@{}}
\toprule\noalign{}
\begin{minipage}[b]{\linewidth}\raggedright
Component
\end{minipage} & \begin{minipage}[b]{\linewidth}\raggedright
Үүрэг
\end{minipage} & \begin{minipage}[b]{\linewidth}\raggedright
Хил ба requirement
\end{minipage} \\
\midrule\noalign{}
\endhead
\bottomrule\noalign{}
\endlastfoot
Plugins/Hooks & Helmet, CORS, rate limit, auth, RBAC, validation &
FR-05, FR-10, NFR-01 \\
Controllers & HTTP input/output, status code, route composition &
FR-01\ldots14 \\
Services & Business rule, state transition, ownership orchestration &
FR-02, FR-09, FR-11\ldots14, NFR-05 \\
Repositories & Supabase query ба DB-specific mapping & FR-01\ldots14,
NFR-01 \\
Shared schemas & Request/response validation contract & FR-03, NFR-03 \\
Audit repository & Admin mutation болон denied-access audit & FR-09,
NFR-01 \\
\end{longtable}

Controller дотор business logic, service дотор шууд database query
хийхгүй. Энэ хил нь ADR-007-ийн гол үр дагавар бөгөөд кодын одоогийн
folder structure-тэй таарна.

\subsection{5.3 Building block → SRS
холбоо}\label{building-block-srs-ux445ux43eux43bux431ux43eux43e}

\begin{longtable}[]{@{}
  >{\raggedright\arraybackslash}p{(\linewidth - 2\tabcolsep) * \real{0.5000}}
  >{\raggedright\arraybackslash}p{(\linewidth - 2\tabcolsep) * \real{0.5000}}@{}}
\toprule\noalign{}
\begin{minipage}[b]{\linewidth}\raggedright
Building block
\end{minipage} & \begin{minipage}[b]{\linewidth}\raggedright
Хариуцах requirement ID
\end{minipage} \\
\midrule\noalign{}
\endhead
\bottomrule\noalign{}
\endlastfoot
Web public/inquiry/history & FR-01, FR-03, FR-04, FR-06, FR-15;
NFR-02/04 \\
Web admin inbox & FR-07, FR-08, FR-09; NFR-01 \\
Web wishlist & FR-10, FR-11, FR-12, FR-13, FR-14; NFR-01 \\
API plugins/auth/validation & FR-02, FR-03, FR-05, FR-10; NFR-01 \\
Inquiry controller/service/repository & FR-01, FR-02, FR-04, FR-06,
FR-07, FR-08, FR-09; NFR-05 planned \\
Cart controller/service/repository & FR-10, FR-11, FR-12, FR-13,
FR-14 \\
PostgreSQL/Auth/audit & FR-01, FR-02, FR-06, FR-09, FR-12; NFR-01 \\
CI/deployment tooling & NFR-03 gap; бүрэлдэхүүнд оноосон боловч
хэрэгжээгүй \\
\end{longtable}

Бүх ID-ийн
\href{https://github.com/Tengis01/Nogoolin/blob/main/docs/requirements/srs-v1.0.md}{SRS
эх} болон
\href{https://github.com/Tengis01/Nogoolin/blob/main/docs/architecture/architecture-traceability.md}{20
мөрийн architecture mapping} холбоотой. Functional mapping coverage
\textbf{15/15 = 100\%}, DoD-ийн ≥80\% буюу 12/15 босгыг хангана.
\textbf{Orphan FR: 0}. NFR-03/05 нь orphan биш, хариуцах building block
байгаа ч хэрэгжилтийн gap; W5 төлөвлөгөөнд үлдээнэ. Coverage нь runtime
test pass гэсэн үг биш.

\section{6. Ажиллах үеийн
зураглал}\label{ux430ux436ux438ux43bux43bux430ux445-ux4afux435ux438ux439ux43d-ux437ux443ux440ux430ux433ux43bux430ux43b}

\subsection{6.1 Inquiry
илгээх}\label{inquiry-ux438ux43bux433ux44dux44dux445}

\begin{figure}[H]
\centering
\includegraphics[width=1\linewidth,height=0.72\textheight,keepaspectratio,alt={C4-D01 --- Inquiry Runtime}]{../tmp/diagrams/c4-d01-inquiry.pdf}
\caption{C4-D01 --- Inquiry Runtime}
\end{figure}

\emph{Mermaid эх:
\href{https://github.com/Tengis01/Nogoolin/blob/main/docs/ICSI438/sem4/diagrams/c4-d01-inquiry.mmd}{c4-d01-inquiry.mmd}.}

Guest эсвэл customer form илгээхэд client ба API ижил Zod contract
ашиглана. Hook нь optional session, input болон rate limit-ийг шалгана.
Service нь inquiry үүсгэж repository-оор PostgreSQL-д хадгална. Invalid
input \texttt{400}, хэтрэлт \texttt{429}, амжилт \texttt{201} байна. Энэ
scenario FR-01\ldots06, NFR-01, NFR-04, NFR-05-тай холбоотой.

\subsection{6.2 Wishlist жагсаалт ба
хадгалалт}\label{wishlist-ux436ux430ux433ux441ux430ux430ux43bux442-ux431ux430-ux445ux430ux434ux433ux430ux43bux430ux43bux442}

\begin{figure}[H]
\centering
\includegraphics[width=1\linewidth,height=0.72\textheight,keepaspectratio,alt={C4-D02 --- Wishlist Runtime}]{../tmp/diagrams/c4-d02-wishlist.pdf}
\caption{C4-D02 --- Wishlist Runtime}
\end{figure}

\emph{Mermaid эх:
\href{https://github.com/Tengis01/Nogoolin/blob/main/docs/ICSI438/sem4/diagrams/c4-d02-wishlist.mmd}{c4-d02-wishlist.mmd}.}

Guest API түвшинд \texttt{401} авна. Customer session баталгаажсаны
дараа service published product болон ownership-ийг шалгаж, repository
unique user-product constraint-тай ажиллана. Remove нь idempotent байна.
Энэ scenario FR-10\ldots14 ба NFR-01-ийг хамарна.

\subsection{6.3 Админ inquiry төлөв
шинэчлэх}\label{ux430ux434ux43cux438ux43d-inquiry-ux442ux4e9ux43bux4e9ux432-ux448ux438ux43dux44dux447ux43bux44dux445}

\begin{figure}[H]
\centering
\includegraphics[width=1\linewidth,height=0.72\textheight,keepaspectratio,alt={C4-D03 --- Admin Inquiry Runtime}]{../tmp/diagrams/c4-d03-admin-inquiry.pdf}
\caption{C4-D03 --- Admin Inquiry Runtime}
\end{figure}

\emph{Mermaid эх:
\href{https://github.com/Tengis01/Nogoolin/blob/main/docs/ICSI438/sem4/diagrams/c4-d03-admin-inquiry.mmd}{c4-d03-admin-inquiry.mmd}.}

Admin inbox command API-д очиход JWT ба role шалгана. Schema
зөвшөөрөгдсөн status утгыг шалгаж, service шинэчлэлт хийж, primary
update амжилттай болсны дараа audit row нэмнэ. Customer \texttt{403},
invalid status \texttt{400}, unknown ID \texttt{404} авна. Энэ scenario
FR-07\ldots09 ба NFR-01-тэй холбоотой.

\section{7. Байршуулалтын
зураглал}\label{ux431ux430ux439ux440ux448ux443ux443ux43bux430ux43bux442ux44bux43d-ux437ux443ux440ux430ux433ux43bux430ux43b}

\begin{longtable}[]{@{}
  >{\raggedright\arraybackslash}p{(\linewidth - 8\tabcolsep) * \real{0.2000}}
  >{\raggedright\arraybackslash}p{(\linewidth - 8\tabcolsep) * \real{0.2000}}
  >{\raggedright\arraybackslash}p{(\linewidth - 8\tabcolsep) * \real{0.2000}}
  >{\raggedright\arraybackslash}p{(\linewidth - 8\tabcolsep) * \real{0.2000}}
  >{\raggedright\arraybackslash}p{(\linewidth - 8\tabcolsep) * \real{0.2000}}@{}}
\toprule\noalign{}
\begin{minipage}[b]{\linewidth}\raggedright
Орчин
\end{minipage} & \begin{minipage}[b]{\linewidth}\raggedright
Web
\end{minipage} & \begin{minipage}[b]{\linewidth}\raggedright
API
\end{minipage} & \begin{minipage}[b]{\linewidth}\raggedright
Data/Auth/Storage
\end{minipage} & \begin{minipage}[b]{\linewidth}\raggedright
Төлөв
\end{minipage} \\
\midrule\noalign{}
\endhead
\bottomrule\noalign{}
\endlastfoot
Local & \texttt{localhost:3000} & \texttt{localhost:3001} & Supabase CLI
Docker stack & Ашиглаж буй development target \\
Production & Vercel & Railway Docker & Supabase cloud, Singapore & Phase
6-д төлөвлөсөн; ажиллаж буй мэт батлаагүй \\
Edge & Cloudflare DNS/WAF & API domain proxy & Хамаарахгүй & Production
target \\
Mobile distribution & Expo dev client / EAS & Ижил REST API & Supabase
Auth & EAS build manual trigger \\
\end{longtable}

Нууц утгууд diagram эсвэл баримтад бичигдэхгүй. Production topology нь
\texttt{docs/phase-0/09-deployment.md}-ийн зорилтот загвар; cloud
project одоогоор байхгүй.

\section{8. Хөндлөн огтлох
ойлголтууд}\label{ux445ux4e9ux43dux434ux43bux4e9ux43d-ux43eux433ux442ux43bux43eux445-ux43eux439ux43bux433ux43eux43bux442ux443ux443ux434}

\begin{longtable}[]{@{}
  >{\raggedright\arraybackslash}p{(\linewidth - 2\tabcolsep) * \real{0.5000}}
  >{\raggedright\arraybackslash}p{(\linewidth - 2\tabcolsep) * \real{0.5000}}@{}}
\toprule\noalign{}
\begin{minipage}[b]{\linewidth}\raggedright
Ойлголт
\end{minipage} & \begin{minipage}[b]{\linewidth}\raggedright
Хэрэгжүүлэх зарчим
\end{minipage} \\
\midrule\noalign{}
\endhead
\bottomrule\noalign{}
\endlastfoot
Authentication & Supabase session; API JWT verification \\
Authorization & \texttt{requireAuth}/\texttt{requireAdmin}, ownership
filter, PostgreSQL RLS \\
Validation & Shared Zod schema; client validation нь UX, server
validation нь заавал \\
Error contract & Тогтвортой HTTP status болон хэрэглэгчид ойлгомжтой
message \\
Audit & Admin mutation ба denied access-ийг append-only log-д үлдээх \\
Configuration & Environment variable; secret repository-д хадгалахгүй \\
Resilience & Reduced motion, WebGL static fallback, idempotent remove \\
Observability & Одоогоор structured application log ба audit; production
monitoring deferred \\
\end{longtable}

\section{9. Архитектурын
шийдвэрүүд}\label{ux430ux440ux445ux438ux442ux435ux43aux442ux443ux440ux44bux43d-ux448ux438ux439ux434ux432ux44dux440ux4afux4afux434}

US-4.3-ийн Tech Stack, Persistence, Interface ангилал бүрд нэг MADR
бичив. ADR-003/005 нь одоогийн шийдвэрийн дэлгэрэнгүй; M4-IF-01 нь
одоогийн REST contract-ийг баримтжуулсан course record. ADR-007 нь
нэмэлт архитектурын үндэслэлээр тусдаа эхэд хадгалагдана.

\begin{longtable}[]{@{}
  >{\raggedright\arraybackslash}p{(\linewidth - 6\tabcolsep) * \real{0.2500}}
  >{\raggedright\arraybackslash}p{(\linewidth - 6\tabcolsep) * \real{0.2500}}
  >{\raggedright\arraybackslash}p{(\linewidth - 6\tabcolsep) * \real{0.2500}}
  >{\raggedright\arraybackslash}p{(\linewidth - 6\tabcolsep) * \real{0.2500}}@{}}
\toprule\noalign{}
\begin{minipage}[b]{\linewidth}\raggedright
ID
\end{minipage} & \begin{minipage}[b]{\linewidth}\raggedright
Шийдвэр
\end{minipage} & \begin{minipage}[b]{\linewidth}\raggedright
Status
\end{minipage} & \begin{minipage}[b]{\linewidth}\raggedright
Requirements/C4 холбоо
\end{minipage} \\
\midrule\noalign{}
\endhead
\bottomrule\noalign{}
\endlastfoot
\href{https://github.com/Tengis01/Nogoolin/blob/main/docs/architecture/adr/adr-003-fastify.md}{ADR-003}
& Fastify over Express & Accepted & FR-01\ldots14, NFR-03/04;
C4-02/03/D01\ldots D03 \\
\href{https://github.com/Tengis01/Nogoolin/blob/main/docs/architecture/adr/adr-005-supabase.md}{ADR-005}
& Supabase over Firebase/self-hosted PostgreSQL & Accepted &
FR-01\ldots14, NFR-01/05; бүх data/auth view \\
\href{https://github.com/Tengis01/Nogoolin/blob/main/docs/architecture/adr/m4-if-01-rest-interface.md}{M4-IF-01}
& Versioned REST + shared Zod contract & Accepted, existing
implementation & FR-01\ldots14, NFR-01/03; C4-02/03 \\
\end{longtable}

ADR-006 Next.js, ADR-010 manual EAS build, ADR-011 local-first DB нь
\href{https://github.com/Tengis01/Nogoolin/blob/main/agent-context/DECISIONS.md}{decision
log}-д хэвээр бөгөөд энэ M4-д бүрэн карт болгон давтаагүй.

\section{10. Чанарын
шаардлага}\label{ux447ux430ux43dux430ux440ux44bux43d-ux448ux430ux430ux440ux434ux43bux430ux433ux430}

\begin{longtable}[]{@{}
  >{\raggedright\arraybackslash}p{(\linewidth - 6\tabcolsep) * \real{0.2500}}
  >{\raggedright\arraybackslash}p{(\linewidth - 6\tabcolsep) * \real{0.2500}}
  >{\raggedright\arraybackslash}p{(\linewidth - 6\tabcolsep) * \real{0.2500}}
  >{\raggedright\arraybackslash}p{(\linewidth - 6\tabcolsep) * \real{0.2500}}@{}}
\toprule\noalign{}
\begin{minipage}[b]{\linewidth}\raggedright
ID
\end{minipage} & \begin{minipage}[b]{\linewidth}\raggedright
Quality scenario
\end{minipage} & \begin{minipage}[b]{\linewidth}\raggedright
Архитектурын хариу
\end{minipage} & \begin{minipage}[b]{\linewidth}\raggedright
Verification төлөв
\end{minipage} \\
\midrule\noalign{}
\endhead
\bottomrule\noalign{}
\endlastfoot
NFR-01 & A хэрэглэгч B-ийн inquiry/wishlist-г харахгүй & JWT + ownership
+ RLS; C4-D01/D02/D03 & Partial inspection; runtime matrix planned \\
NFR-02 & Reduced motion үед CTA 1 секундэд бэлэн & Web container
static/reduced-motion path & Implemented; test planned \\
NFR-03 & DB integration suite skip бол CI fail & Build/deployment gate;
local Supabase dependency & Gap; planned W9 \\
NFR-04 & 10 run-ийн 9-д каталог ≤2.0 s & SSR/SSG, Fastify, indexed
PostgreSQL query & Defined; measurement not executed \\
NFR-05 & Ижил idempotency key 24 цагт нэг inquiry & Service/repository
key store шаардлагатай & Planned; not implemented \\
\end{longtable}

\section{11. Эрсдэл ба техникийн
өр}\label{ux44dux440ux441ux434ux44dux43b-ux431ux430-ux442ux435ux445ux43dux438ux43aux438ux439ux43d-ux4e9ux440}

\begin{longtable}[]{@{}
  >{\raggedright\arraybackslash}p{(\linewidth - 4\tabcolsep) * \real{0.3333}}
  >{\raggedright\arraybackslash}p{(\linewidth - 4\tabcolsep) * \real{0.3333}}
  >{\raggedright\arraybackslash}p{(\linewidth - 4\tabcolsep) * \real{0.3333}}@{}}
\toprule\noalign{}
\begin{minipage}[b]{\linewidth}\raggedright
Эрсдэл/өр
\end{minipage} & \begin{minipage}[b]{\linewidth}\raggedright
Нөлөө
\end{minipage} & \begin{minipage}[b]{\linewidth}\raggedright
Одоогийн хариу
\end{minipage} \\
\midrule\noalign{}
\endhead
\bottomrule\noalign{}
\endlastfoot
CI local DB provision хийхгүй & Integration test бүгд skip хийгээд green
болох & NFR-03 gap гэж ил тод тэмдэглэсэн \\
Inquiry idempotency key store байхгүй & Retry давхар inquiry үүсгэж
магадгүй & NFR-05 planned; ADR шаардлагатай байж болно \\
Phase 0-ийн Express/хуучин path & Architecture drift, буруу diagram &
Код ба decision log-ийг M4 source болгосон \\
Production service үүсээгүй & Deployment зураг нотолгоо мэт ойлгогдох &
§7-д target гэж тэмдэглэсэн \\
Mobile inquiry/wishlist дутуу & Container capability зөрөх & C4-02 болон
status-д gap гэж тэмдэглэсэн \\
Mermaid renderer/layout өөрчлөгдөх & Renderer update layout эвдэх & CLI
11.17.0 pin, lockfile, SVG diff ба targeted visual QA \\
\end{longtable}

\section{12. Нэр
томьёо}\label{ux43dux44dux440-ux442ux43eux43cux44cux451ux43e}

\begin{longtable}[]{@{}
  >{\raggedright\arraybackslash}p{(\linewidth - 2\tabcolsep) * \real{0.5000}}
  >{\raggedright\arraybackslash}p{(\linewidth - 2\tabcolsep) * \real{0.5000}}@{}}
\toprule\noalign{}
\begin{minipage}[b]{\linewidth}\raggedright
Нэр
\end{minipage} & \begin{minipage}[b]{\linewidth}\raggedright
Тайлбар
\end{minipage} \\
\midrule\noalign{}
\endhead
\bottomrule\noalign{}
\endlastfoot
arc42 & Architecture documentation-ийн 12 хэсэгтэй бүтэц \\
C4 & System Context, Container, Component, Code гэсэн zoom түвшний
загвар \\
Container & Тусдаа ажиллах эсвэл deploy хийх application/data store \\
Component & Container доторх тодорхой үүрэгтэй нэгж \\
Dynamic view & Нэг feature ажиллах үеийн элементүүдийн дараалал \\
ADR & Нэг архитектурын шийдвэрийн context, choice, consequence
бүртгэл \\
RLS & PostgreSQL Row Level Security \\
Layered monolith & Нэг deployable API дотор хариуцлагаар тусгаарласан
давхаргууд \\
\end{longtable}

\subsection{Эх
сурвалж}\label{ux44dux445-ux441ux443ux440ux432ux430ux43bux436}

\begin{itemize}
\tightlist
\item
  \href{https://arc42.org/download/}{arc42 template v9 download}
\item
  \href{https://docs.arc42.org/section-3/}{arc42 §3 --- Context and
  Scope}
\item
  \href{https://docs.arc42.org/section-5/}{arc42 §5 --- Building Block
  View}
\item
  \href{https://docs.arc42.org/section-6/}{arc42 §6 --- Runtime View}
\item
  \href{https://c4model.com/diagrams/system-context}{C4 model --- System
  Context}
\item
  \href{https://c4model.com/diagrams/dynamic}{C4 model --- Dynamic
  diagram}
\item
  \href{https://mermaid.js.org/syntax/c4}{Mermaid C4 syntax}
\item
  \href{https://adr.github.io/adr-templates/}{ADR templates}
\end{itemize}

\newpage

\section{Architecture Traceability --- M4 v0.1
draft}\label{architecture-traceability-m4-v0.1-draft}

Энэ хүснэгт нь
\href{https://github.com/Tengis01/Nogoolin/blob/main/docs/requirements/traceability-matrix-v1.0.md}{Seminar
3-ын 20 мөртэй traceability matrix}-ийн ID-г өөрчлөлгүй архитектурын
element, arc42 section, C4 view болон ADR-тай холбоно. Энэ нь
verification execution-ийн шинэ нотолгоо биш.

\begingroup\footnotesize\setlength{\tabcolsep}{2pt}\sloppy

\begin{longtable}[]{@{}
  >{\raggedright\arraybackslash}p{(\linewidth - 10\tabcolsep) * \real{0.1667}}
  >{\raggedright\arraybackslash}p{(\linewidth - 10\tabcolsep) * \real{0.1667}}
  >{\raggedright\arraybackslash}p{(\linewidth - 10\tabcolsep) * \real{0.1667}}
  >{\raggedright\arraybackslash}p{(\linewidth - 10\tabcolsep) * \real{0.1667}}
  >{\raggedright\arraybackslash}p{(\linewidth - 10\tabcolsep) * \real{0.1667}}
  >{\raggedright\arraybackslash}p{(\linewidth - 10\tabcolsep) * \real{0.1667}}@{}}
\toprule\noalign{}
\begin{minipage}[b]{\linewidth}\raggedright
ID
\end{minipage} & \begin{minipage}[b]{\linewidth}\raggedright
Architecture responsibility
\end{minipage} & \begin{minipage}[b]{\linewidth}\raggedright
arc42
\end{minipage} & \begin{minipage}[b]{\linewidth}\raggedright
C4 view
\end{minipage} & \begin{minipage}[b]{\linewidth}\raggedright
ADR
\end{minipage} & \begin{minipage}[b]{\linewidth}\raggedright
Architecture status
\end{minipage} \\
\midrule\noalign{}
\endhead
\bottomrule\noalign{}
\endlastfoot
\href{https://github.com/Tengis01/Nogoolin/blob/main/docs/requirements/srs-v1.0.md\#fr-01--зочны-inquiry-үүсгэх}{FR-01}
& Inquiry Controller → Service → Repository & §5, §6.1 & C4-02, C4-03,
C4-D01 & 003, 005, 007 & Implemented path; test planned \\
\href{https://github.com/Tengis01/Nogoolin/blob/main/docs/requirements/srs-v1.0.md\#fr-02--нэвтэрсэн-хэрэглэгчийн-inquiry-г-холбох}{FR-02}
& optionalAuth, session-derived customer ID & §5.2, §6.1, §8 & C4-03,
C4-D01 & 003, 005, 007 & Implemented; test planned \\
\href{https://github.com/Tengis01/Nogoolin/blob/main/docs/requirements/srs-v1.0.md\#fr-03--inquiry-form-ийг-баталгаажуулах}{FR-03}
& Shared Zod schema and validation hook & §5.2, §8 & C4-03, C4-D01 &
003, 007 & Implemented; test planned \\
\href{https://github.com/Tengis01/Nogoolin/blob/main/docs/requirements/srs-v1.0.md\#fr-04--inquiry-баталгааг-харуулах}{FR-04}
& Inquiry response and Web success state & §5.1, §6.1 & C4-02, C4-D01 &
003, 007 & Implemented; test planned \\
\href{https://github.com/Tengis01/Nogoolin/blob/main/docs/requirements/srs-v1.0.md\#fr-05--inquiry-rate-limit-хэрэгжүүлэх}{FR-05}
& Fastify route rate limit & §5.2, §6.1, §8 & C4-03, C4-D01 & 003, 007 &
Implemented; test planned \\
\href{https://github.com/Tengis01/Nogoolin/blob/main/docs/requirements/srs-v1.0.md\#fr-06--өөрийн-inquiry-history-г-харах}{FR-06}
& Authenticated own-inquiry query and RLS & §5.2, §6.1, §8 & C4-03,
C4-D01 & 005, 007 & Implemented; test planned \\
\href{https://github.com/Tengis01/Nogoolin/blob/main/docs/requirements/srs-v1.0.md\#fr-07--админ-inquiry-inbox-харах}{FR-07}
& Admin inbox controller/repository & §5.1, §6.3 & C4-02, C4-03, C4-D03
& 003, 005, 007 & Implemented; test planned \\
\href{https://github.com/Tengis01/Nogoolin/blob/main/docs/requirements/srs-v1.0.md\#fr-08--админ-inquiry-шүүх}{FR-08}
& Admin query validation and filtering & §5.2, §6.3 & C4-03, C4-D03 &
003, 007 & Implemented; test planned \\
\href{https://github.com/Tengis01/Nogoolin/blob/main/docs/requirements/srs-v1.0.md\#fr-09--inquiry-төлөв-шинэчлэх}{FR-09}
& Status service and audit repository & §5.2, §6.3, §8 & C4-03, C4-D03 &
005, 007 & Implemented; test planned \\
\href{https://github.com/Tengis01/Nogoolin/blob/main/docs/requirements/srs-v1.0.md\#fr-10--wishlist-д-нэвтрэлт-шаардах}{FR-10}
& Web route guard and API requireAuth & §5.1, §6.2, §8 & C4-02, C4-03,
C4-D02 & 003, 005, 007 & Implemented; test planned \\
\href{https://github.com/Tengis01/Nogoolin/blob/main/docs/requirements/srs-v1.0.md\#fr-11--нийтлэгдсэн-бүтээгдэхүүн-хадгалах}{FR-11}
& Published-product rule in cart service & §5.2, §6.2 & C4-03, C4-D02 &
005, 007 & Implemented; test planned \\
\href{https://github.com/Tengis01/Nogoolin/blob/main/docs/requirements/srs-v1.0.md\#fr-12--wishlist-давхардлыг-хориглох}{FR-12}
& Service/repository plus unique DB pair & §5.2, §6.2 & C4-03, C4-D02 &
005, 007 & Implemented; test planned \\
\href{https://github.com/Tengis01/Nogoolin/blob/main/docs/requirements/srs-v1.0.md\#fr-13--wishlist-жагсаах}{FR-13}
& Customer-scoped cart list and Web grid & §5.1, §6.2 & C4-02, C4-03,
C4-D02 & 003, 005, 007 & Implemented; test planned \\
\href{https://github.com/Tengis01/Nogoolin/blob/main/docs/requirements/srs-v1.0.md\#fr-14--wishlist-ээс-хасах}{FR-14}
& Idempotent remove service/repository & §5.2, §6.2 & C4-03, C4-D02 &
005, 007 & Implemented; test planned \\
\href{https://github.com/Tengis01/Nogoolin/blob/main/docs/requirements/srs-v1.0.md\#fr-15--hero-fallback-аас-каталогт-хүрэх}{FR-15}
& Web intro config and static fallback & §5.1, §8, §10 & C4-02 &
Existing ADR-006 & Partial; test planned \\
\href{https://github.com/Tengis01/Nogoolin/blob/main/docs/requirements/srs-v1.0.md\#nfr-01--хувийн-өгөгдлийг-тусгаарлах}{NFR-01}
& JWT, RBAC, ownership, RLS defense-in-depth & §3, §5, §6, §8, §10 &
C4-01\ldots03, C4-D01\ldots03 & 003, 005, 007 & Partial inspection \\
\href{https://github.com/Tengis01/Nogoolin/blob/main/docs/requirements/srs-v1.0.md\#nfr-02--хөдөлгөөн-багасгах-тохиргоо}{NFR-02}
& Reduced-motion and static Web path & §5.1, §8, §10 & C4-02 & Existing
ADR-006 & Implemented; test planned \\
\href{https://github.com/Tengis01/Nogoolin/blob/main/docs/requirements/srs-v1.0.md\#nfr-03--deploy-оос-өмнөх-тестийн-хаалт}{NFR-03}
& Local DB integration and negative CI gate & §7, §10, §11 & C4-02 &
Existing ADR-010/011 & Gap; planned W9 \\
\href{https://github.com/Tengis01/Nogoolin/blob/main/docs/requirements/srs-v1.0.md\#nfr-04--каталогийн-хуудсын-ачааллын-хугацаа}{NFR-04}
& SSR/Web, Fastify API and indexed data path & §4, §5, §10 & C4-02,
C4-03 & 003, 005, 007 & Defined; not measured \\
\href{https://github.com/Tengis01/Nogoolin/blob/main/docs/requirements/srs-v1.0.md\#nfr-05--inquiry-retry-давхардал-үүсгэхгүй-байх}{NFR-05}
& Inquiry service idempotency boundary & §5.2, §6.1, §10, §11 & C4-03,
C4-D01 & 005, 007 & Planned; not implemented \\
\end{longtable}

\endgroup

\subsection{Coverage summary}\label{coverage-summary}

\begin{itemize}
\tightlist
\item
  Requirement rows: \textbf{20 / 20}.
\item
  Functional building-block coverage: \textbf{15/15 = 100\%}; required
  ≥80\% = 12/15.
\item
  Orphan FR: \textbf{0}. NFR-03 ба NFR-05 implementation gap-ийг W5-д
  үргэлжлүүлнэ.
\item
  Static views: C4-01, C4-02, C4-03.
\item
  Runtime views: C4-D01, C4-D02, C4-D03.
\item
  Main MADRs: ADR-003 (Tech Stack), ADR-005 (Persistence), M4-IF-01
  (Interface). ADR-007 нь supporting decision.
\item
  \texttt{implemented} нь code path байгааг, \texttt{verified} нь
  ажиллуулсан нотолгоог илэрхийлнэ; энэ хүснэгт шинэ requirement-ийг
  verified гэж зарлаагүй.
\end{itemize}

\newpage

\section{ADR-003 --- Fastify-г API framework болгон
сонгох}\label{adr-003-fastify-ux433-api-framework-ux431ux43eux43bux433ux43eux43d-ux441ux43eux43dux433ux43eux445}

\begin{itemize}
\tightlist
\item
  \textbf{Status:} Accepted
\item
  \textbf{Date:} 2026-06; M4 expanded 2026-09-28
\item
  \textbf{Decision owner:} Тэнгис
\item
  \textbf{Related requirements:} FR-01\ldots14, NFR-03, NFR-04
\item
  \textbf{Related views:} C4-02, C4-03, C4-D01\ldots D03
\end{itemize}

\subsection{Context}\label{context}

Nogoolin-д web, admin, mobile гурвыг үйлчлэх TypeScript REST API
хэрэгтэй. Solo developer учраас framework нь бага boilerplate-тай,
schema validation, security hook, plugin encapsulation болон integration
test-д тохиромжтой байх шаардлагатай.

\subsection{Decision Drivers}\label{decision-drivers}

Нэг хөгжүүлэгчийн ажиллагаа, шалгаж болох interface, өгөгдөл хамгаалалт
болон requirements coverage.

\subsection{Considered options}\label{considered-options}

\begin{longtable}[]{@{}
  >{\raggedright\arraybackslash}p{(\linewidth - 4\tabcolsep) * \real{0.3333}}
  >{\raggedright\arraybackslash}p{(\linewidth - 4\tabcolsep) * \real{0.3333}}
  >{\raggedright\arraybackslash}p{(\linewidth - 4\tabcolsep) * \real{0.3333}}@{}}
\toprule\noalign{}
\begin{minipage}[b]{\linewidth}\raggedright
Сонголт
\end{minipage} & \begin{minipage}[b]{\linewidth}\raggedright
Давуу тал
\end{minipage} & \begin{minipage}[b]{\linewidth}\raggedright
Сул тал
\end{minipage} \\
\midrule\noalign{}
\endhead
\bottomrule\noalign{}
\endlastfoot
Fastify & TypeScript support, plugin/hook model, хурдан request path &
Express-ээс жижиг ecosystem; plugin lifecycle сурах шаардлагатай \\
Express & Өргөн ecosystem, олон жишээ & Validation/security
composition-г гараар их зохион байгуулна \\
Next.js route handlers only & Web-тэй нэг deploy & Mobile/admin API хил,
repository/service structure бүдгэрнэ \\
\end{longtable}

\subsection{Decision}\label{decision}

Fastify + TypeScript-ийг port 3001 дээрх \texttt{/api/v1/} REST API
framework болгоно. Helmet, CORS, rate-limit, auth, validation-ийг
plugin/hook-д; HTTP mapping-ийг controller-д; business rule-ийг
service-д; data access-ийг repository-д байрлуулна.

\subsection{Consequences}\label{consequences}

\begin{itemize}
\tightlist
\item
  Web, admin, mobile нэг тогтвортой API contract ашиглана.
\item
  Security болон validation route handler бүрт давтагдахгүй.
\item
  Plugin order, encapsulation болон error mapping-ийг test-ээр хамгаалах
  шаардлагатай.
\item
  Phase 0 дахь Express нэршил хуучирсан; шинэ architecture view
  Fastify-г ашиглана.
\end{itemize}

\subsection{Persona pain point ба
confirmation}\label{persona-pain-point-ux431ux430-confirmation}

\href{https://github.com/Tengis01/Nogoolin/blob/main/docs/ICSI438/sem1/wiki/02_audience_persona.md}{W1
persona}-ийн \textbf{P3: local ба pipeline-ийн ялгааг оношлох; P-NG-03}
хэрэгцээтэй холбов. Энэ нь RAG кейсээс Nogoolin руу шилжүүлсэн
тайлбарласан холбоо; уг алдаа Nogoolin-д болсон гэж үзээгүй.
\href{https://github.com/Tengis01/Nogoolin/blob/main/docs/ICSI438/sem2/wiki/persona-and-pivot.md}{W2
pivot} шилжилтийг тайлбарлана.

Confirmation: decision log болон одоогийн code structure-тай static
check хийсэн; шинэ runtime test энэ ажлаар ажиллуулаагүй.

\newpage

\section{ADR-005 --- Supabase-г data/auth/storage platform болгон
сонгох}\label{adr-005-supabase-ux433-dataauthstorage-platform-ux431ux43eux43bux433ux43eux43d-ux441ux43eux43dux433ux43eux445}

\begin{itemize}
\tightlist
\item
  \textbf{Status:} Accepted
\item
  \textbf{Date:} 2026-07 (Phase 0); M4 expanded 2026-09-28
\item
  \textbf{Decision owner:} Тэнгис
\item
  \textbf{Related requirements:} FR-01\ldots14, NFR-01, NFR-05
\item
  \textbf{Related views:} C4-01, C4-02, C4-03, C4-D01\ldots D03
\end{itemize}

\subsection{Context}\label{context-1}

Nogoolin-д relational product/inquiry/cart өгөгдөл, user authentication,
media storage болон row-level authorization хэрэгтэй. Solo developer
тусдаа database, auth server, object storage ажиллуулахад operational
ачаалал өндөр.

\subsection{Decision Drivers}\label{decision-drivers-1}

Нэг хөгжүүлэгчийн ажиллагаа, шалгаж болох interface, өгөгдөл хамгаалалт
болон requirements coverage.

\subsection{Considered options}\label{considered-options-1}

\begin{longtable}[]{@{}
  >{\raggedright\arraybackslash}p{(\linewidth - 4\tabcolsep) * \real{0.3333}}
  >{\raggedright\arraybackslash}p{(\linewidth - 4\tabcolsep) * \real{0.3333}}
  >{\raggedright\arraybackslash}p{(\linewidth - 4\tabcolsep) * \real{0.3333}}@{}}
\toprule\noalign{}
\begin{minipage}[b]{\linewidth}\raggedright
Сонголт
\end{minipage} & \begin{minipage}[b]{\linewidth}\raggedright
Давуу тал
\end{minipage} & \begin{minipage}[b]{\linewidth}\raggedright
Сул тал
\end{minipage} \\
\midrule\noalign{}
\endhead
\bottomrule\noalign{}
\endlastfoot
Supabase & PostgreSQL, Auth, Storage, RLS нэг platform; local CLI stack
& Platform-specific client/API; cloud dependency \\
Firebase & Managed auth/data, mobile ecosystem & Relational query болон
SQL/RLS загварт тохиромж бага \\
Self-hosted PostgreSQL + auth/storage & Бүрэн хяналт & Setup, patch,
backup, auth, storage ажиллагаа их \\
\end{longtable}

\subsection{Decision}\label{decision-1}

Supabase PostgreSQL, Auth, Storage-г ашиглана. Local development ба
security test-ийг Supabase CLI stack дээр хийнэ; cloud project Phase 6
хүртэл deferred. API repository interface нь Supabase-specific query-г
service-ээс тусгаарлана.

\subsection{Consequences}\label{consequences-1}

\begin{itemize}
\tightlist
\item
  NFR-01-д JWT, ownership болон RLS-ийг хамтад нь хэрэгжүүлэх боломжтой.
\item
  Migration append-only байх ба local reset дээр дахин ажиллах ёстой.
\item
  \texttt{service\_role} key зөвхөн server environment-д байна.
\item
  Provider солиход repository implementation болон auth integration
  өөрчлөгдөнө; business service-ийг аль болох хэвээр үлдээнэ.
\end{itemize}

\subsection{Persona pain point ба
confirmation}\label{persona-pain-point-ux431ux430-confirmation-1}

\href{https://github.com/Tengis01/Nogoolin/blob/main/docs/ICSI438/sem1/wiki/02_audience_persona.md}{W1
persona}-ийн \textbf{P1: session алдагдах; P-NG-01, мөн P2:
schema/producer contract нийцэл} хэрэгцээтэй холбов. Энэ нь RAG кейсээс
Nogoolin руу шилжүүлсэн тайлбарласан холбоо; уг алдаа Nogoolin-д болсон
гэж үзээгүй.
\href{https://github.com/Tengis01/Nogoolin/blob/main/docs/ICSI438/sem2/wiki/persona-and-pivot.md}{W2
pivot} шилжилтийг тайлбарлана.

Confirmation: decision log болон одоогийн code structure-тай static
check хийсэн; шинэ runtime test энэ ажлаар ажиллуулаагүй.

\newpage

\section{M4-IF-01 --- REST interface ба shared validation
contract}\label{m4-if-01-rest-interface-ux431ux430-shared-validation-contract}

\begin{itemize}
\tightlist
\item
  \textbf{Status:} Accepted --- existing implementation documented
\item
  \textbf{Date:} 2026-09-28
\item
  \textbf{Decision owner:} Тэнгис
\item
  \textbf{Category:} Interface; MADR course record
\item
  \textbf{Requirements:} FR-01\ldots14, NFR-01, NFR-03
\item
  \textbf{Views:} C4-02, C4-03, C4-D01\ldots03
\end{itemize}

\subsection{Context}\label{context-2}

Web, admin, mobile нэг API ашиглана. Field, HTTP status болон auth
boundary зөрвөл client зөв request илгээсэн ч өөр орчинд буруу үр дүн
авна. W1-ийн
\href{https://github.com/Tengis01/Nogoolin/blob/main/docs/ICSI438/sem1/wiki/02_audience_persona.md}{P2
schema/producer нийцэл, P3 local/pipeline ялгаа}-г шинэ төсөлд contract
drift-ийн эрсдэлтэй холбов. W1 vector dimension нь Nogoolin-ийн
өгөгдлийн талбар биш.

\subsection{Decision Drivers}\label{decision-drivers-2}

Client бүрт ойлгомжтой HTTP interface, reusable TypeScript schema,
testable auth/status boundary, нэг хөгжүүлэгчид бага tooling хэрэгтэй.

\subsection{Considered Options}\label{considered-options-2}

\begin{enumerate}
\def\labelenumi{\arabic{enumi}.}
\tightlist
\item
  Versioned REST JSON + shared Zod schema.
\item
  GraphQL schema/server нэмж client query composition хийх.
\item
  Client бүр тусдаа request/validation model хадгалах.
\end{enumerate}

\subsection{Decision}\label{decision-2}

Одоогийн \texttt{/api/v1/} REST JSON interface болон shared Zod
package-г баримтжуулна. Inquiry POST \texttt{201}, invalid input
\texttt{400}, guest auth scope \texttt{401}, non-admin \texttt{403},
unknown product \texttt{404}, rate limit \texttt{429} байна. Wishlist
\texttt{/cart} нэртэй боловч checkout биш. HTTP API нь runtime
interface; Zod package нь library тул C4 Container болгон зурахгүй.

\subsection{Consequences}\label{consequences-2}

\begin{itemize}
\tightlist
\item
  Web/API contract нэг schema-тай, алдааны төлөв тодорхой болно.
\item
  Schema шинэчлэхэд consumer bundler ба test-ийг хамт шалгана.
\item
  Mobile одоогоор зарим type-only import хэрэглэж байгаа; runtime
  validation parity-г бүрэн болсон гэж зарлахгүй.
\item
  Legacy OpenAPI-г W5-д одоогийн inquiry/cart contract-тай нийцүүлэх
  шаардлагатай.
\end{itemize}

\subsection{Confirmation}\label{confirmation}

Одоогийн controllers, shared schemas, web API client-ийг статик байдлаар
тулгав. Энэ record шинэ product decision, шинэ endpoint эсвэл runtime
test үүсгэхгүй.

\newpage

\section{UE-4 --- Container зураглалын Before/After
засвар}\label{ue-4-container-ux437ux443ux440ux430ux433ux43bux430ux43bux44bux43d-beforeafter-ux437ux430ux441ux432ux430ux440}

\subsection{Аргачлал}\label{ux430ux440ux433ux430ux447ux43bux430ux43b}

Bhatti Ch.6-ийн гурван алхмыг хэрэглэв: эхлээд элемент, хилийн ноорог
гаргах; уншигч хаанаас эхлэхийг тодруулах; box/arrow бүрийг утгатай
label-тай болгох. Handout нь х.112/113/116 гэж заасан боловч локал номын
хэвлэмэл хуудас дээр эдгээр гурван дэд гарчиг х.116-д хамт байна. Энд
цаасан дээр зурсан эсвэл хүнээр review хийлгэсэн нотолгоо зохиогоогүй;
editable Mermaid эхийг ноорог болгон ашиглав.

\subsection{Before --- зориуд оруулсан
anti-pattern}\label{before-ux437ux43eux440ux438ux443ux434-ux43eux440ux443ux443ux43bux441ux430ux43d-anti-pattern}

\begin{figure}[H]
\centering
\includegraphics[width=1\linewidth,height=0.72\textheight,keepaspectratio,alt={UE-4 Before --- Magic Container}]{../tmp/diagrams/ue4-before-container.pdf}
\caption{UE-4 Before --- Magic Container}
\end{figure}

``Magic Container'' нь web, mobile, API-ийн ялгаа, технологи,
хариуцлагыг нууж байна. ``Does stuff'' arrow нь ямар мэдээлэл дамжуулж
байгааг хэлэхгүй. Энэ нь зөвхөн UE-4 сургалтын алдаатай зураг; ажиллаж
буй system design гэж үзэхгүй.

\subsection{After --- нэр, технологи, хариуцлага тодорхой
болсон}\label{after-ux43dux44dux440-ux442ux435ux445ux43dux43eux43bux43eux433ux438-ux445ux430ux440ux438ux443ux446ux43bux430ux433ux430-ux442ux43eux434ux43eux440ux445ux43eux439-ux431ux43eux43bux441ux43eux43d}

\begin{figure}[H]
\centering
\includegraphics[width=1\linewidth,height=0.72\textheight,keepaspectratio,alt={UE-4 After --- Nogoolin Container View}]{../tmp/diagrams/c4-02-container.pdf}
\caption{UE-4 After --- Nogoolin Container View}
\end{figure}

Chinchilla Ch.2 х.22--23-ын architecture reference нь харилцан
хамааралтай хэсгүүдийг тайлбарлах зорилготой. After зурагт Web/Admin,
Mobile, REST API-г тусдаа deployable unit болгон задалж, Supabase
services-ийг хөрш системээр ялгав. Actor, request, authentication болон
persistence замын label тодорхой болсон. Shared schema нь library учраас
container биш; §5.2 component view-д reference хийнэ.

\begin{longtable}[]{@{}
  >{\raggedright\arraybackslash}p{(\linewidth - 4\tabcolsep) * \real{0.3333}}
  >{\raggedright\arraybackslash}p{(\linewidth - 4\tabcolsep) * \real{0.3333}}
  >{\raggedright\arraybackslash}p{(\linewidth - 4\tabcolsep) * \real{0.3333}}@{}}
\toprule\noalign{}
\begin{minipage}[b]{\linewidth}\raggedright
Өмнөх зөрүү
\end{minipage} & \begin{minipage}[b]{\linewidth}\raggedright
Засвар
\end{minipage} & \begin{minipage}[b]{\linewidth}\raggedright
Үр дүн
\end{minipage} \\
\midrule\noalign{}
\endhead
\bottomrule\noalign{}
\endlastfoot
Нэг ``Magic'' box & Web, Mobile, API болгож салгасан & Deployable
хариуцлага харагдана \\
``Various technologies'' & Next.js, Expo, Fastify & Stack тодорхой \\
``Does stuff'' & JSON API, session, query/media label & Data flow
ойлгомжтой \\
Бүх боломж бэлэн мэт & Mobile inquiry/wishlist incomplete & Бодит
capability-г ялгасан \\
\end{longtable}

\newpage

\section{UE4-ADR-01 --- Magic Container зураглалыг
засах}\label{ue4-adr-01-magic-container-ux437ux443ux440ux430ux433ux43bux430ux43bux44bux433-ux437ux430ux441ux430ux445}

\textbf{Status:} Accepted for exercise · \textbf{Date:} 2026-09-28 ·
\textbf{Owner:} Тэнгис

\subsection{Context}\label{context-3}

Before зураг нэг ``Magic Container''-д бүх frontend/backend хариуцлагыг
багтаасан. Хөгжүүлэгч алдаа гарвал аль deployable unit-ийг шалгахаа
мэдэхгүй. Энэ нь W1-ийн
\href{https://github.com/Tengis01/Nogoolin/blob/main/docs/ICSI438/sem1/wiki/02_audience_persona.md}{P3
local/pipeline ялгааг оношлох} хэрэгцээтэй холбоотой.

\subsection{Decision Drivers}\label{decision-drivers-3}

Зураг унших эхлэл, deployable boundary, name/technology/responsibility
болон data flow-г тодорхой болгох. Bhatti Ch.6 х.116-ын label зарчим,
Chinchilla Ch.2 х.22--23-ын architecture reference зорилгыг хэрэглэв.

\subsection{Considered Options}\label{considered-options-3}

\begin{enumerate}
\def\labelenumi{\arabic{enumi}.}
\tightlist
\item
  Magic box-д урт тайлбар нэмэх.
\item
  Web/Admin, Mobile, REST API-ийг тусдаа container болгон задалж хөрш
  platform services-ийг ялгах.
\end{enumerate}

\subsection{Decision}\label{decision-3}

Хоёрдугаар сонголтыг сонгов. After зураг нь гурван Nogoolin application
container-тай; Supabase platform нь Auth, PostgreSQL, Storage хөрш
services-ээр харагдана. Shared Zod schema-г library гэж тайлбарлаж,
deployable container болгож дүрслэхгүй. Actor болон arrow бүр утгатай
label-тай байна.

\subsection{Consequences}\label{consequences-3}

Уншигч Web → API → Data урсгал болон session замыг ялгана. Mobile-ийн
дутуу capability-г ил тод харуулна. Зураг засах нь product кодын
refactoring эсвэл runtime validation биш; deployment boundary-г
өөрчлөөгүй.

\subsection{Confirmation}\label{confirmation-1}

Before/After Mermaid эх, SVG болон §5 building block хүснэгтийг тулгав.
Nogoolin applications \textbf{3 ≤ 8}. FR mapping \textbf{15/15} бөгөөд
NFR-03/05 gap-ийг W5-д үлдээв.

\newpage

\section{Seminar 4 reflection}\label{seminar-4-reflection}

\subsection{1. Bhatti-ийн ямар ойлголт миний дадлыг өөрчилсөн
бэ?}\label{bhatti-ux438ux439ux43d-ux44fux43cux430ux440-ux43eux439ux43bux433ux43eux43bux442-ux43cux438ux43dux438ux439-ux434ux430ux434ux43bux44bux433-ux4e9ux4e9ux440ux447ux438ux43bux441ux4e9ux43d-ux431ux44d}

Bhatti Ch.6 х.116-ын ``Use labels'' зарчим хамгийн их нөлөөлөв. Би өмнө
нь нэртэй хайрцаг ба сум байхад архитектур ойлгомжтой гэж үздэг байсан.
UE-4-ийн Magic Container дээр нэр, технологи, хариуцлага болон сумны
data flow байхгүй үед зураг мэдээлэл өгдөггүйг харлаа. After хувилбарт
``Uses'', ``Does stuff''-ийг JSON request, session, persistence гэсэн
бодит харилцаагаар сольж, уншигчийн эхлэх actor-ийг тодорхой болгов.
Номын х.117 SVG зөвлөмжийн дагуу зургаа scalable vector хэлбэрээр
хадгалав.

\subsection{2. Chinchilla-ийн ямар ойлголтыг W16-аас хойш ашиглах
вэ?}\label{chinchilla-ux438ux439ux43d-ux44fux43cux430ux440-ux43eux439ux43bux433ux43eux43bux442ux44bux433-w16-ux430ux430ux441-ux445ux43eux439ux448-ux430ux448ux438ux433ux43bux430ux445-ux432ux44d}

Chinchilla Ch.4 х.47--51-ын navigation ба hierarchy зарчмыг
үргэлжлүүлнэ. Architecture-ийг нэг урт тайлбарт нуухын оронд arc42
§3-аас §5 руу, requirement-ээс building block, ADR руу холбоосоор шилжих
боломжтой болгов. Ch.2 х.22--23-ын дагуу architecture reference нь
хэсгүүдийн харилцаа, сонгосон шийдлийн үндэслэлийг тайлбарлах ёстой;
эхлэх tutorial-ийг орлохгүй. Дараагийн төсөлд мөн diagram, decision,
requirement холбоосыг хамт хадгална.

\subsection{3. Ном ба бодит ажиллагааны хамгийн том зөрүү хаана байна
вэ?}\label{ux43dux43eux43c-ux431ux430-ux431ux43eux434ux438ux442-ux430ux436ux438ux43bux43bux430ux433ux430ux430ux43dux44b-ux445ux430ux43cux433ux438ux439ux43d-ux442ux43eux43c-ux437ux4e9ux440ux4afux4af-ux445ux430ux430ux43dux430-ux431ux430ux439ux43dux430-ux432ux44d}

Ном зураглалаа уншигчаар шалгах, label-ийг ойлгодог эсэхийг туршихыг
зөвлөдөг. Одоогоор би зураг, код, requirement-ээ өөрөө тулгасан; гаднын
reviewer-ийн санал байхгүй. Мөн 100\% FR mapping нь код ажиллаж, тест
өнгөрсөн гэсэн баталгаа биш. NFR-03 CI gate, NFR-05 idempotency gap
хэвээр. Иймээс M4-д coverage ба runtime evidence-ийг ялгаж, W5-д
interface drift болон эдгээр gap-ийг үргэлжлүүлэхээр тэмдэглэв.

\subsection{Standup ба retrospective --- өөрийн
тэмдэглэл}\label{standup-ux431ux430-retrospective-ux4e9ux4e9ux440ux438ux439ux43d-ux442ux44dux43cux434ux44dux433ux43bux44dux43b}

Өнөөдөр §1/3/5-ыг зааварт нийцүүлж, зургаа Mermaid эхтэй болголоо.
Үндсэн гурван ADR-ийн scope тодорхой; нэмэлт тодруулга нь W1-ийн ганц
persona-ийг гурван бизнес role-той холбосон mapping-ийг багш хэрхэн
үнэлэх тухай байна. Хамгийн их тус болсон хэсэг нь §5: хариуцлага ба FR
mapping нэг дор харагдана. Дараагийн sprint-д ADR-ийн confirmation болон
interface drift шалгалтыг илүү тодорхой болгоно. Энэ нь багийн хурал
болсон тухай протокол биш.

\end{document}
